\subsection{Qubits}
Como se introduce en~\cite{diracBraket}, los cúbits se escriben normalmente en la llamada notación bra-ket.
Cada cúbit se representa como un ket, de la forma $\ket{a}$, que denota un vector en $\Q$.

Por ejemplo, los estados de la base computacional se escriben como $\ket{0}$ y $\ket{1}$, llamados ket 0 y ket 1, respectivamente (véase~\cite{NielsenChuang}).
Así, cualquier cúbit puede describirse como una combinación lineal de estos dos estados de la base: $\ket{a} =\alpha_0\ket{0} +\alpha_1 \ket{1}$, donde $\alpha_0,  \alpha_1 \in \mathbb{C}$ con $ |\alpha_0|^{2}+|\alpha_1|^{2}=1$.

\begin{notation}
  Para cualquier cúbit $\ket{a}$, se usa la siguiente notación:
  \[a_{0}=\braket{0}{a}\text{ y }a_{1}=\braket{1}{a},\]
\end{notation}
por lo que el cúbit $\ket{a}$ puede escribirse como:
\begin{equation}
  \label{eq:cúbit-coordenadas}
  \ket{a}=a_0\ket{0}+a_1\ket{1}
\end{equation}

Para otras bases, como por ejemplo $\set{\ket{+}, \ket{-}}$, donde $\ket{+}=\frac{\ket{0}+\ket{1}}{\sqrt{2}}$ y $\ket{-}=\frac{\ket{0}-\ket{1}}{\sqrt{2}}$, un cúbit $\ket{a}$ puede describirse como:
\begin{equation*}
  \ket{a}=a_+\ket{+}+a_-\ket{-}
\end{equation*}
